\noindent
{\Large \textbf{\section{Математическое описание}}} \par

bf16 - формат хранения вещественного числа в 16 битах: 1 бит знака $s$, 8 бит смещённой экспоненты $E$ (смещение 127) и 7 бит мантиссы $M$. Число хранится в нормализованном виде $1.M \cdot 2^{E-127}$, единица перед точкой не хранится явно. Вещественное число вычисляется по формуле:
\[
    x = (-1)^s \cdot (1 + M) \cdot 2^{E-127}.
\]
Крайние значения экспоненты зарезервированы: $E=0$ при $M=0$ даёт ноль, иначе - денормализованное число $x=(-1)^s \cdot M \cdot 2^{1-127}$; $E=255$ даёт $\pm\infty$ (при $M=0$) или NaN (при $M\neq0$).

\vspace{6pt}
\textbf{Перевод из десятичной формы в bf16.} Целая часть числа переводится в двоичный вид делением на 2, дробная - умножением на 2 (на каждом шаге дробная часть умножается на 2, целая часть результата - очередной бит, остаток идёт в следующий шаг). Полученная запись нормализуется к виду $1.xxxxx \cdot 2^E$ сдвигом до первой единицы, после чего 7 бит мантиссы округляются по следующему за ними биту (округление к ближайшему), а $E$ смещается на 127.

\vspace{6pt}
\textbf{Перевод из bf16 в десятичную форму.} Обратная операция: 16 бит разбираются на знак, экспоненту и мантиссу, после чего значение восстанавливается по формуле выше (с учётом особых случаев нуля, денормализованных чисел, бесконечности и NaN).

\vspace{6pt}
Перевод целых чисел (без дробной части) выполняется классическим способом - делением на 2 с накоплением остатков (в прямом направлении) и суммированием степеней двойки (в обратном).

\newpage